\subsection{加速运动中的固有时累积}
匀速公式用于整段旅程之前，要检查钟是否全程相对所选惯性系以同一速率运动，即 $\gamma$ 是否恒定。速度随时间变化时，每个足够短的时间段满足
\[
 \dd\tau=\sqrt{1-v(t)^2/c^2}\,\dd t,\qquad
 \tau=\int_{t_A}^{t_B}\sqrt{1-v(t)^2/c^2}\,\dd t.
\]
这给出理想钟沿既定世界线的累计读数；不是把全程平均速度放入 $\gamma$。只要各段之间钟的读数连续，便可逐段相加。
例如地面与飞船在A相遇，飞船以 $0.60c$ 匀速飞出3.00年（地面计时），在远端短时掉头，再以同速返回并在B重逢。忽略掉头所占时间，地面钟累计6.00年，飞船钟累计
\[
 \tau=2(3.00\,\mathrm{yr})\sqrt{1-0.60^2}=4.80\,\mathrm{yr}.
\]
A、B是同一对重逢事件，两只钟的累计读数可以直接比较。飞船去、回两段对应不同惯性系，不能把“运动钟慢”的互易说法在中途无条件拼接；差异应由各自世界线上的固有时积分计算。实际有限时长掉头还需把掉头段的积分加入。
